\section{Algebra-Based Definition of RML}%
\label{sec:expressiveness}

This section
	introduces
	%defines
an algorithm to convert mappings defined using the mapping language RML (specifically, version~1.1.2~\cite{Dimou2024:RMLSpec}) into our
	mapping
algebra.
	%By providing this algorithm, we show
	%With this algorithm, we show
	This algorithm shows
that \emph{our algebra is at least as expressive as RML}. Moreover, through this algorithm, in combination with the algebra, we provide
	a
	%the first
\emph{formal definition of the semantics of RML mappings}.
	%To the best of our knowledge, this is the first work that captures the semantics of RML formally.
	%Regarding the latter, to the best of our knowledge, the semantics of RML has not been captured formally before.
%
	Hereafter, we describe the algorithm step by step, while a complete pseudocode representation is given as Algorithm~\ref{alg:rml2algebra}.
	%A complete pseudocode representation of the algorithm is given as Algorithm~\ref{alg:rml2algebra}, which we describe in the following, step by step.
For the description, we assume
	%the reader is familiar
	familiarity
with the concepts~of~RML%
	~\cite{dimou_ldow_2014,Dimou2024:RMLSpec,IglesiasMolina2023:RMLOntology}%
.

\begin{algorithm}[t!]
\begin{algorithm}[t]
    \SetArgSty{textnormal}
    \DontPrintSemicolon
	\caption{Translates an RML mapping into our mapping algebra.}
	\label{alg:rml2algebra}
        \KwData{%
            $\symRMLGraph$ - an RDF graph (assumed to describe an RML mapping) \;
            \hspace{10mm}
            $\LTB$ - an injective function
            	that maps
            	%mapping
            every string to a unique blank~node
        }
        \KwResult{a mapping relation }

		$\omega_\mathsf{spog} \gets$ initially empty algebraic expression \;

		\ForEach{$\triplesMapIRI_\mathsf{tm}\! \in \symAllIRIs \cup \symAllBNodes$ for which
			%there is a triple
		$(\triplesMapIRI_\mathsf{tm}, \ttlSmaller{rdf:type}, \ttlSmaller{rr:TriplesMap}) \in \symRMLGraph$}
        {

			$\omega \gets \sourceOp{\dataSource}{\queryExpr}{\symAttrQueryMap}$\!,
				where
				%with
			$(\dataSource, \queryExpr,\symAttrQueryMap) =$ \textsc{SrcOpParams}($\triplesMapIRI_\mathsf{tm},\symRMLGraph$)%
            \tcp*{Algorithm~\ref{alg:rml_source_op}}

			$u_\textsf{sm} \gets o$, where $o \in \symAllIRIs \cup \symAllBNodes$ such that $(\triplesMapIRI_\mathsf{tm},\ttlSmaller{rr:subjectMap},o) \in \symRMLGraph$ \;

			$\omega \gets \extOpAttr{\subjAttr\!}(\omega)$,
				%where
				with
			$\extExpr \!=\! \textsc{CreateExtExpr}(u_\textsf{sm}, \symRMLGraph, \LTB, \symAttrQueryMap)$ 
            \tcp*[r]{Algorithm~\ref{alg:rml_termmap_extend_expr}} 

			$u_\textsf{pom} \gets o$, where $o \in \symAllIRIs \cup \symAllBNodes$ such that $(\triplesMapIRI_\mathsf{tm},\ttlSmaller{rr:predicateObjectMap},o) \in \symRMLGraph$ \; 

			$u_\textsf{pm} \gets o$, where $o \in \symAllIRIs \cup \symAllBNodes$ such that $(u_\textsf{pom},\ttlSmaller{rr:predicateMap},o) \in \symRMLGraph$ \;

			$\omega \gets \extOpAttr{\predAttr\!}(\omega)$,
				%where
				with
			$\extExpr \!=\! \textsc{CreateExtExpr}(u_\textsf{pm}, \symRMLGraph, \LTB, \symAttrQueryMap)$
            \tcp*[r]{Algorithm~\ref{alg:rml_termmap_extend_expr}} 

            $u_\mathsf{om} \gets o$, where $o \in \symAllIRIs \cup \symAllBNodes$ such that $(u_\textsf{pom},\ttlSmaller{rr:objectMap},o) \in \symRMLGraph$ \;

			\If{there is a $u_\mathsf{ptm} \in \symAllIRIs \cup \symAllBNodes$ such that $(u_\mathsf{om},\ttlSmaller{rr:parentTriplesMap},u_\mathsf{ptm}) \in \symRMLGraph$}
            {

				$\symJoinAttrPairs \gets \emptyset$ \;
				\ForEach{$u_\mathsf{jc} \in \symAllIRIs \cup \symAllBNodes$ for which
					%there is a triple
				$(u_\mathsf{om},\ttlSmaller{rr:joinCondition},u_\mathsf{jc}) \in \symRMLGraph$}
                {

					$\attr \gets \symAttrQueryMap^{-1}(\lex)$, where $(u_\mathsf{jc},\ttlSmaller{rr:child},o) \in \symRMLGraph$ with $o = (\lex,\dt) \in \symAllLiterals$ \;

					$\attr'\! \gets \symAttrQueryMap^{-1}(\lex)$, where $(u_\mathsf{jc},\ttlSmaller{rr:parent},o) \in \symRMLGraph$ with $o = (\lex,\dt) \in \symAllLiterals$ \;

					$\symJoinAttrPairs \gets \symJoinAttrPairs \cup \{ (\attr, \attr') \}$ \;
                }

				$\omega \gets \fctEquiJoinOpBig{\omega}{ \sourceOp{\dataSource'\!}{\queryExpr'\!}{\symAttrQueryMap'} }$,
					%where
					with
				$(\dataSource'\!, \queryExpr'\!,\symAttrQueryMap') \!=\!$ \textsc{SrcOpParams}($u_\mathsf{ptm},\symRMLGraph$)
               \;% \tcp*[r]{Algorithm~\ref{alg:rml_source_op}} 

				$u_\mathsf{sptm} \gets o$, where $o \in \symAllIRIs \cup \symAllBNodes$ such that $(u_\mathsf{ptm},\ttlSmaller{rr:subjectMap},o) \in \symRMLGraph$ \;

				$\omega \gets \extOpAttr{\objAttr\!}(\omega)$, where $\extExpr = \textsc{CreateExtExpr}(u_\mathsf{sptm}, \symRMLGraph, \LTB, \symAttrQueryMap')$ \;
            }
            \Else{
				$\omega \gets \extOpAttr{\objAttr\!}(\omega)$, where $\extExpr = \textsc{CreateExtExpr}(u_\textsf{om}, \symRMLGraph, \LTB, \symAttrQueryMap)$ \;
            }


			\If{there is a $u_\mathsf{gm} \in \symAllIRIs \cup \symAllBNodes$ such that
            $(u_\textsf{sm},\ttlSmaller{rr:graphMap},u_\mathsf{gm}) \in \symRMLGraph$ or \endgraf\hspace{49mm} $(u_\textsf{pom},\ttlSmaller{rr:graphMap},u_\mathsf{gm}) \in \symRMLGraph$}
            {

				$\omega \gets \extOpAttr{\graphAttr\!}(\omega)$, where $\extExpr = \textsc{CreateExtExpr}(u_\mathsf{gm}, \symRMLGraph, \LTB, \symAttrQueryMap)$ \;
            }
            \Else{
				$\omega \gets \extOpAttr{\graphAttr\!}(\omega)$, where $\extExpr$ is the IRI \ttlSmaller{rr:defaultGraph} \;
            }

			$\omega_\mathsf{spog} \gets \fctUnionBig{\omega_\mathsf{spog}}{ \fctProjectOp{ \{\subjAttr,\predAttr,\objAttr,\graphAttr\} }{\omega} }$ \;
        }
        \Return{$\omega_\mathsf{spog}$}
\end{algorithm}

\end{algorithm}

\smallskip\noindent\textbf{\itshape Input and Output.}
Since RML mappings are captured as RDF graphs that use
	%a dedicated vocabulary (the so-called RML ontology~\cite{IglesiasMolina2023:RMLOntology})%
	the RML ontology~\cite{IglesiasMolina2023:RMLOntology}%
, the main input to our algorithm is such an RDF graph. As a second input, we assume an injective function~%
	$\LTB \!: \symAllStrings \rightarrow \symAllBNodes$
	%$\LTB$
that maps every string to a unique blank node (which shall become relevant later). The output
	%of the algorithm
is a mapping relation that is obtained by
	applying
operators of our algebra and that can be converted into an RDF dataset as per Definition~\ref{def:generated_dataset}.


\smallskip\noindent\textbf{\itshape Normalization Phase.}
To minimize the
	%cases and variations of RML
	variations of RML descriptions
%
	%that the algorithm has to consider, it
	to be considered, the algorithm
begins by converting the given RML mapping into a \emph{normal form} in which i)~no shortcut properties
	%for class IRIs and for constant-valued term maps
are used, ii)~%
	%every referencing object map has
	all referencing object maps have
a join condition, and iii)~every triples map has only a single predicate-object map with a single predicate map, a single object map, and at most one graph map%
	%~(which may be associated either with the predicate-object map itself or with its object map)%
.
Every RML mapping can be converted into this normal form by applying the following sequence of rewriting steps, which we have adopted from Rodríguez-Muro and Rezk~\cite{Mariano2015SPARQL2SQL}, with adaptations specific to RML.
\PaperVersion{While we present these steps informally,
Appendix~\ref{appdx:queries} captures them formally in terms of update~queries.}
\ExtendedVersion{We present these steps with the accompanying update~queries which captures the steps formally.}
None of these steps changes the meaning of the defined~mapping.

\begin{enumerate}
	\item
	\label{item:NormalizationStep:ClassExpansion}
	Every class IRI definition is expanded into a predicate-object map.
        \ExtendedVersion{(Listing~\ref{lst:sparql_class_expand})}

	\item
	\label{item:NormalizationStep:ConstantsExpansion}
	All shortcut properties for constant-valued term maps are expanded.
        \ExtendedVersion{(Listing~\ref{lst:sparql_shortcut})}

	\item
	\label{item:NormalizationStep:POMDuplication}
	Predicate-object maps with multiple predicate maps or multiple object maps are duplicated to have a single predicate map and a single object map each.
        \ExtendedVersion{(Listing~\ref{lst:sparql_deduplication})}

	\item
	\label{item:NormalizationStep:JoinRemoval}
	Referencing object maps without a join condition are replaced by an object map with the subject IRI of the parent triples map as the reference value.
        \ExtendedVersion{(Listing~\ref{lst:sparql_selfjoin})}

	\item
	\label{item:NormalizationStep:TMDuplication1}
	Every triples map with multiple predicate-object maps is duplicated such that each of the resulting triples maps has a single predicate-object map~only.
        \ExtendedVersion{(Listing~\ref{lst:sparql_tm_singlepom})}

	\item
	\label{item:NormalizationStep:TMDuplication2}
	Every triples map with multiple graph maps is duplicated such that each of the resulting triples maps has one of these graph maps.
        \ExtendedVersion{(Listing~\ref{lst:sparql_graphmap_norm})}
\end{enumerate}



\begin{figure*}[t]
        \centering
	\begin{minipage}[t]{0.43\textwidth}
		\centering
		\begin{lstlisting}[language=sparql, label={lst:sparql_class_expand}, caption={Normalization step~\ref{item:NormalizationStep:ClassExpansion} (expand shortcuts for class IRIs%
			%\ into predicate-object maps%
		).}]
DELETE { ?sm rr:class ?sm_class . }
INSERT { ?tm rr:predicateObjectMap [
      rr:predicateMap [
          rr:constant rdf:type ;
          rr:termType rr:IRI ; ] ;
      rr:objectMap [
         rr:termType rr:IRI ;
         rr:constant ?sm_class ; 
    ] ; ] . }
WHERE { ?tm rr:subjectMap ?sm . 
        ?sm rr:class ?sm_class . }
\end{lstlisting}
	\end{minipage}%
    \hfill%
    \begin{minipage}[t]{0.53\textwidth}
        \centering
        \begin{lstlisting}[language=sparql, caption={Normalization step~\ref{item:NormalizationStep:POMDuplication} (duplicate multi-predicate-object maps
	%\ into singleton predicate-object maps%
	\ into singletons%
).},
label={lst:sparql_deduplication}]
DELETE{ ?tm rr:predicateObjectMap  ?pom. 
  ?pom rr:predicateMap ?pm;
  rr:objectMap ?om; rr:graphMap ?gm. }
INSERT{ ?tm rr:predicateObjectMap [
      rr:predicateMap ?pm;
      rr:objectMap ?om; 
      rr:graphMap ?gm; ]. }
WHERE{ ?tm rr:predicateObjectMap  ?pom. 
  ?pom rr:predicateMap ?pm. 
  ?pom rr:objectMap ?om. 
  OPTIONAL {?pom rr:graphMap ?gm.} } 
\end{lstlisting}
    \end{minipage}
%
    \begin{minipage}[t]{\textwidth}
\vspace{-2mm}
\begin{lstlisting}[language=sparql,  caption={Normalization step~\ref{item:NormalizationStep:ConstantsExpansion} (expand shortcuts
for constant-valued term maps).},
label={lst:sparql_shortcut}]
DELETE {  ?tm   rr:subject ?sm_constant .   ?pompm rr:predicate ?pm_constant .
          ?termMap rr:graph ?gm_constant .  ?pomom rr:object  ?om_constant .   }
INSERT { ?tm      rr:subjectMap   [ rr:constant ?sm_constant; ].
         ?pompm   rr:predicateMap [ rr:constant ?pm_constant; ].
         ?pomom   rr:objectMap    [ rr:constant ?om_constant; ].
         ?termMap rr:graphMap     [ rr:constant ?gm_constant; ]. }
WHERE { { ?tm      rr:subject   ?sm_constant }
  UNION { ?pompm   rr:predicate ?pm_constant }
  UNION { ?pomom   rr:object    ?om_constant }
  UNION { ?termMap rr:graph     ?gm_constant }  }
\end{lstlisting}
    \end{minipage}
\vspace{-2mm}
    \begin{minipage}[t]{0.4\textwidth}
\begin{lstlisting}[language=sparql, caption={Normalization step~\ref{item:NormalizationStep:JoinRemoval} (replace referencing object maps
  that do not have a join condition%
  %without join conditions%
).},
label={lst:sparql_selfjoin}]
DELETE{ ?om rr:parentTriplesMap ?ptm .}
INSERT{ ?om rr:reference  ?ref ;
    rr:template ?template; 
    rr:constant ?const ;
    rr:termType rr:IRI . }
WHERE{ ?om rr:parentTriplesMap ?ptm .
?ptm rr:subjectMap ?sm . 
OPTIONAL{?sm rr:reference ?ref .}
OPTIONAL{?sm rr:template ?template .}
OPTIONAL{?sm rr:constant ?const .}
FILTER NOT EXISTS 
{?om rr:joinCondition ?jc .} } 
\end{lstlisting}
    \end{minipage}
    \hfill
    \begin{minipage}[t]{0.5\textwidth}
\begin{lstlisting}[language=sparql, caption={Normalization step~\ref{item:NormalizationStep:TMDuplication1} (duplicate triples maps that contain multiple 
pred\-i\-cate-ob\-ject maps).},
label={lst:sparql_tm_singlepom}]
DELETE{ ?tm a rr:TriplesMap;
    rml:logicalSource ?ls; rr:subjectMap ?sm;
    rr:predicateObjectMap ?pom . }
INSERT{ [] a rr:TriplesMap;
   rml:logicalSource ?ls;
   rr:subjectMap ?sm; 
   rr:predicateObjectMap ?pom. }
WHERE{ ?tm rml:logicalSource ?ls;
    rr:subjectMap ?sm;
    rr:predicateObjectMap ?pom . }
\end{lstlisting}
    \end{minipage}
\end{figure*}

\begin{lstlisting}[language=sparql, float=t, caption={Normalization step~\ref{item:NormalizationStep:TMDuplication2} (duplicate triples maps
   %that contain multiple graph~maps%
   %containing multiple graph~maps%
   to contain one graph map%
 ).},
label={lst:sparql_graphmap_norm}]
DELETE { ?tm rml:logicalSource ?ls ; rr:subjectMap ?sm ; rr:predicateObjectMap ?pom ;
            a rrTriplesMap .
         ?sm   rr:graphMap ?gm. 
         ?pom  rr:graphMap ?gm2 . }
INSERT { [] rml:logicalSource ?ls ; rr:subjectMap ?sm ; rr:predicateObjectMap ?pom ;
            a rr:TriplesMap .
         ?sm rr:graphMap ?gm . 
         [] rml:logicalSource ?ls ; rr:subjectMap ?sm ; rr:predicateObjectMap ?pom ;
            a rr:TriplesMap .
         ?sm rr:graphMap ?gm2 . }
WHERE { ?tm rml:logicalSource ?ls ; rr:subjectMap ?sm ; rr:predicateObjectMap ?pom .
        ?sm   rr:graphMap ?gm . 
        ?pom  rr:graphMap ?gm2 . }
\end{lstlisting}

% \smallskip
\noindent\textbf{\itshape Main Loop.}
After normalization, the algorithm iterates over all triples maps defined by the (normalized) RML mapping (lines~\ref{line:rml2algebra:main_loop_begin}--\ref{line:rml2algebra:main_loop_end}). For each of them, it
	%defines a mapping relation based on algebra operators (as described below) 
	combines algebra operators, as described below, to define a mapping relation that, in the end, has the attributes $\subjAttr$, $\predAttr$, $\objAttr$, and $\graphAttr$.
These mapping relations are combined using union (line~\ref{line:rml2algebra:union}), resulting in a mapping relation (line~\ref{line:rml2algebra:return})
	that captures the complete RDF dataset defined by the given RML mapping.
	%that, by applying Definition~\ref{def:generated_dataset}, captures all RDF triples produced by the RML mapping for the logical sources.


\smallskip\noindent\textbf{\itshape Translation of the Logical Source.}
As the first step for each triples map, the algorithm creates a source operator
	%based on the logical source defined for the triples map~%
(line~\ref{line:rml2algebra:source}). The first two parameters for this operator---the data source~$\dataSource$ and the query expression $\queryExpr$ for selecting relevant context objects (cf.\ Definition~\ref{def:source_operator})---depend on the logical source defined for the triples map. Since RML is extensible regarding the types of logical sources,
	%we cannot give a complete formal definition of this aspect of our translation algorithm. Instead, we
	we have to
abstract the creation of these two parameters by the following function.

\begin{definition}
	\label{def:initSrcAndRootQuery}
	\normalfont
	Given an IRI or blank node $\triplesMapIRI_\mathsf{tm}$ (assumed to denote an RML triples map) and an RDF graph~$\symRMLGraph$ (assumed to describe an RML mapping, including
		%the triples map~%
	$\triplesMapIRI_\mathsf{tm}$),
	$\textsc{SrcAndRootQuery}(\triplesMapIRI_\mathsf{tm},\symRMLGraph)$ is a pair $(\dataSource, \queryExpr)$ of a data source $\dataSource = \dataSourceTuple$ with $\sourceType = \sourceTypeTuple$ and a query expression~$\queryExpr \in \symDataAcc$, as extracted from the description of the logical source of $\triplesMapIRI_\mathsf{tm}$~in~$\symRMLGraph$.
\end{definition}

\counterwithout{lstlisting}{chapter}
\begin{lstlisting}[%
	language=rdf,
	caption={RML mapping with
		%a single
		one
	triples map (written in Turtle, prefix decl.\ omitted).},
	label={lst:rml_example},
	float=t%
]
ex:tm rml:logicalSource [ rml:source "data.csv";
                          rml:referenceFormulation ql:CSV ];
      rr:subjectMap [ rr:template "http://example.com/person_{ID}";
                      rr:termType rr:IRI ];
      rr:predicateObjectMap [ rr:predicate rdfs:label;
                              rr:objectMap [rml:reference "Name"] ].
\end{lstlisting}

\begin{example}
	\label{ex:initSrcAndRootQuery}
		%Consider
		Given
	an RDF graph $\symRMLGraph_\mathsf{ex}$ that
		%describes an RML mapping including
		contains
	the
		%following triples.
		%\begin{align*}
		%	& \quad (\ttl{ex:tm}, \, \ttl{rml:logicalSource}, \, \ttl{ex:ls}) \in \symRMLGraph_\mathsf{ex}
		%	\\
		%	&(\ttl{ex:ls}, \, \ttl{rml:source}, \, \ttl{"data.csv"}) \in \symRMLGraph_\mathsf{ex}
		%	\\
		%	&(\ttl{ex:ls}, \, \ttl{rml:referenceFormulation}, \, \ttl{ql:CSV}) \in \symRMLGraph_\mathsf{ex}
		%\end{align*}
		%Then%
		RML mapping of Listing~\ref{lst:rml_example}%
	, $\textsc{SrcAndRootQuery}(\ttl{ex:tm},\symRMLGraph_\mathsf{ex}) = (\dataSource_\mathsf{ex}', \varepsilon)$ where $\dataSource_\mathsf{ex}' = (\sourceTypeX{\mathsf{csv}},\bigDataObject_\mathsf{ex}')$ such that $\sourceTypeX{\mathsf{csv}}$ is the source type for CSV-based data sources (cf.\ Example~\ref{ex:source_type_csv}), $\bigDataObject_\mathsf{ex}'$ is the CSV file mentioned in the second triple, and $\varepsilon$ is the empty string.
\end{example}

The third parameter for the source operator---the
	partial
function~$\symAttrQueryMap$
	that maps attributes to query expressions
	%from attributes to query expressions
(cf.\ Definition~\ref{def:source_operator})---is populated by
	%the function \textsc{ExtractQueries}, defined in Algorithm~\ref{alg:query_extraction}. This function
	Algorithm~\ref{alg:query_extraction}, which
extracts the query expressions
	from the term maps of
	%for
the given triples map. In particular, the term maps may use query expressions as references (lines~\ref{line:query_extraction:from_references_begin}--\ref{line:query_extraction:from_references_end} of Algorithm~\ref{alg:query_extraction}), in
	string
templates (lines~\ref{line:query_extraction:from_templates_begin}--\ref{line:query_extraction:from_templates_end}%
	%\ of Algorithm~\ref{alg:query_extraction}%
), and
	as child and parent queries of
	%in
join conditions (lines~\ref{line:query_extraction:from_joins_begin}--\ref{line:query_extraction:from_joins_end}%
	%\ of Algorithm~\ref{alg:query_extraction}%
).
	To focus only on the term maps of the given triples map, the extraction is restricted to the relevant
	%To focus only on the relevant term maps, the extraction is restricted to the relevant
	%We emphasize that this extraction focuses only on a triples-map-specific
subgraph of the given RDF graph (line~\ref{line:query_extraction:subgraph} of Algorithm~\ref{alg:query_extraction})%
	%. Formally, we define such subgraphs as follows.
	, which we define formally~as~follows.

\begin{algorithm}[t]
\begin{algorithm}[htbp]
    \SetArgSty{textnormal}
    \DontPrintSemicolon
	\caption{\textsc{ExtractQueries}: Extract query expressions from triples maps.}
	\label{alg:query_extraction}
        \KwData{
            $\triplesMapIRI_\textsf{tm}$ - IRI or blank node of the triples map
	            %for which query expressions are to be extracted
	            to extract query expressions from \;
		    $\symRMLGraph$ - an RDF graph (assumed to contain a \emph{normalized} RML description)
        }

        \KwResult{a set of query expressions}

		$\symRMLGraph_\mathsf{tm} \gets$ the $\triplesMapIRI_\mathsf{tm}$-rooted subgraph of $\symRMLGraph$ \tcp*{Definition~\ref{def:rooted_subgraph}}

        $Q \gets \emptyset$ \tcp*{Initially empty set of query expressions}

		\ForAll{$s \in \symAllIRIs \cup \symAllBNodes$ and $o \in \symAllLiterals$ for which there is a triple $(s, \ttlSmaller{rml:reference},o) \in \symRMLGraph_\mathsf{tm}$}
        {
		    $Q \gets Q \cup \{ \lex \}$, where $o = (\lex,\dt)$
        }
		\ForAll{$s \in \symAllIRIs \cup \symAllBNodes$ and $o \in \symAllLiterals$ for which there is a triple $(s, \ttlSmaller{rr:template}, o) \in \symRMLGraph_\mathsf{tm}$}
        {
			\State $Q \gets Q \cup Q'\!$, where $o = (\lex,\dt)$ and $Q'$
				is the set of
				%contains
			all substrings of $\lex$ that are enclosed by "\{" and "\}"
        }
		
		$T_\textrm{c} \gets \{ (s,p,o) \in \symRMLGraph_\mathsf{tm} \, | \, p \text{ is the IRI } \ttlSmaller{rr:child} \text{ and } o \in \symAllLiterals \}$ \;

		$T_\textrm{p} \gets \{(s,p,o) \in \symRMLGraph \, | \, p \text{ is the IRI } \ttlSmaller{rr:parent}, o \in \symAllLiterals, \text{ and there is an } s'\! \in \symAllIRIs\cup\symAllBNodes \text{ such}$ $\text{that } (s'\!, \ttlSmaller{rr:joinCondition}, s) \in \symRMLGraph \text{ and } (s'\!, \ttlSmaller{rr:parentTriplesMap}, \triplesMapIRI_\textsf{tm}) \in \symRMLGraph \}$ \;

        \ForAll{$(s,p,o) \in (T_\textrm{c} \cup T_\textrm{p})$}
        {
		    $Q \gets Q \cup \{ \lex \}$, where $o = (\lex,\dt)$
        }
        \Return{$Q$}
\end{algorithm}


\end{algorithm}

\begin{definition}
	\label{def:rooted_subgraph}
	\normalfont
	Let $\symRMLGraph$ be an RDF graph and
		%$\triplesMapIRI$ be an IRI or a blank node.
		$\triplesMapIRI \in \symAllIRIs \cup \symAllBNodes$.
	The \textbf{$\triplesMapIRI$-rooted subgraph of $\symRMLGraph$}
	is an RDF graph $\symRMLGraph'\! \subseteq \symRMLGraph$
	that is defined recursively as follows:
	\begin{enumerate}
		\item $\symRMLGraph'$ contains every triple $(s,p,o) \in \symRMLGraph$ for which it holds that $s = \triplesMapIRI$.
		\item $\symRMLGraph'$ contains every triple $(s,p,o) \in \symRMLGraph$ for which there
			%exists a
			%is a
			%already exists a
			is already another
		triple $(s'\!,p'\!,o') \in \symRMLGraph'$ such that $s = o'$ and $p'$ is not the IRI \ttl{rr:parentTriplesMap}.
	\end{enumerate}
\end{definition}

\begin{example}
	\label{ex:ExtractQueries}
	For the RDF graph $\symRMLGraph_\mathsf{ex}$ that contains the RML mapping of Listing~\ref{lst:rml_example},
	$\textsc{ExtractQueries}(\ttl{ex:tm},\symRMLGraph_\mathsf{ex})$ returns $\symAttrQueryMap = \{ \attr_1 \rightarrow \textsf{\small ID}, \attr_2 \rightarrow \textsf{\small Name} \}$,
	where $\attr_1$ and $\attr_2$ are arbitrary attributes such that $\attr_1 \neq \attr_2$ and $\attr_1, \attr_2 \notin \{ \subjAttr, \predAttr, \objAttr, \graphAttr \}$.
\end{example}


% The source operator constructed for the triples map in this step creates a mapping relation with values obtained via the extracted query expressions (along the lines of Example~\ref{ex:source_operator}). This relation can now be extended by using these values to create the RDF terms defined by the term maps of the triples map.
% \smallskip
\noindent\textbf{\itshape Translation of the Subject Map and the Predicate Map.}
The source operator constructed
	%for the triples map
in the previous step creates a mapping relation with values obtained via the extracted query expressions (along the lines of Example~\ref{ex:source_operator}).
	This relation can now be extended by using these values
	%These values can then be used
to create RDF terms as defined by the term maps of the triples map. To this end, for the
	subject~(term) map and the predicate~(term)
	%subject map and the predicate
map, the algorithm adds an extend operator, respectively~(lines~\ref{line:rml2algebra:extend_for_subject}--\ref{line:rml2algebra:extend_for_predicate} in Algorithm~\ref{alg:rml2algebra}). The construction of the extend expressions of these operators
	(see Definitions~\ref{def:ExtendExpressions:Syntax} and~\ref{def:extend_operator})
is captured in Algorithm~\ref{alg:rml_termmap_extend_expr} and uses the following extension functions\PaperVersion{, which are defined formally in Appendix~\ref{appdx:functions}.}\ExtendedVersion{.}

\ExtendedVersion{
For every $v \in (\symTermUniverse \cup \error)$, $\dt \in (\symTermUniverse \cup \error)$, and $v'\! \in (\symTermUniverse \cup \error)$, we define:
%
\begin{equation*}
	\fctToIRI{v} =
	\begin{cases}
		\mathit{lex}
		& \text{if $v$ is a literal $\literalTuple$ such that $\lex$ is a valid IRI}
		\\[-1mm]
		& \text{and $\dt$ is the IRI \ttl{xsd:string},}
		\\
		\error
		& \text{else.}
	\end{cases}
\end{equation*}
\begin{equation*}
	\fctToBNode{v} =
	\begin{cases}
		\LTB(\lex)
		& \text{if $v$ is a literal $\literalTuple$
%such that
s.t.\
%
%$\dt$ is the IRI \ttl{xsd:string}%
$\dt = \ttl{xsd:string}$%
        ,}
        \\ 
		\error
		& \text{else.}
	\end{cases}
\end{equation*}
\begin{equation*}
    \fctToLiteral{v, \dt} = 
    \begin{cases}
        (\lex, \dt) &\text{if $v$ is a literal~$(\lex, \dt')$
%such that
s.t.\
%
%$\dt$ is the IRI \ttl{xsd:string}%
$\dt = \ttl{xsd:string}$%
        }, \\
        \error & \text{else.}
    \end{cases}
\end{equation*}
\begin{equation*}
    \fctConcat{v}{v'} = 
    \begin{cases}
        (\lex \circ \lex'\!, \ttl{xsd:string})  &\text{if $v$ and $v'$ are literals $(\lex,\dt)$ and}
		\\[-1mm]
        & \text{$(\lex'\!,\dt')$, respectively,
          %such that both}
          such that}
		\\[-1mm]
        &
          %\text{$\dt$ and $\dt'$ are the IRI \ttl{xsd:string},}
          \dt = \dt'\! = \ttl{xsd:string},
        \\ 
        \error &\text{else,}
		\\
    \end{cases}
\end{equation*}
where $\lex \circ \lex'$ is the string obtained from concatenating the strings $\lex$ and~$\lex'\!$.
}

\PaperVersion{
\begin{itemize}
	\item
	$\toIRI$
		%is a unary extension function that
	converts string literals representing IRIs into these IRIs. 

	\item
	$\toBNode$%
		%\ is a unary extension function for which we assume a function $\LTB: \symAllStrings \rightarrow \symAllBNodes$ that is injective. Then, $\toBNode$ uses $\LTB$ to convert string literals into blank nodes.
		%\ converts string literals into blank nodes, for which we assume the aforementioned injective function $\LTB: \symAllStrings \rightarrow \symAllBNodes$ (see the discussion of the input of Algorithm~\ref{alg:rml2algebra}).
		%, parameterized by the aforementioned injective function $\LTB$ (see the discussion of the input of Algorithm~\ref{alg:rml2algebra}), converts string literals into blank nodes.
		, parameterized by the aforementioned function $\LTB$ (see the discussion of the input of Algorithm~\ref{alg:rml2algebra}), converts string literals into blank~nodes.

	\item
	$\toLiteral$ is a binary extension function that converts string literals into literals with a datatype IRI that is given as the second argument.

	\item
	$\concat$ is a binary extension function that concatenates two string literals. 
\end{itemize}
}

\begin{algorithm}[t]
\begin{algorithm}[h!tbp]
	\caption{\textsc{CreateExtExpr}:
		%Creates an extend expression for an RML term~map.}
		Creates an extend expression for a term~map.}
		%Creates extend expressions for term~maps.}
	\label{alg:rml_termmap_extend_expr}
	\begin{algorithmic}[1]
		\Require 
            $u$ - IRI or blank node of the term map
	            %for which the extend expression is to be created
	            to create the extend expression for
            \Statex \qquad $\symRMLGraph$ - an RDF graph (assumed to describe an RML mapping, incl.\ term map~$u$)
            \Statex \qquad $\LTB$ - an injective function that maps all literals to blank nodes
            \Statex \qquad $\symAttrQueryMap$ - an injective partial function that maps attributes to query expressions
		\Ensure an extend expression (as per Definition~\ref{def:ExtendExpressions:Syntax})
        \State $\extExpr \gets $ initially empty extend expression

        \If{$\symRMLGraph$ contains a triple $(u, \ttlSmaller{rr:reference}, o)$ such that $o$ is a literal $\literalTuple$}
    		\State $\extExpr \gets a$, where $a$ is the attribute in $\fctDom{\symAttrQueryMap}$ such that $\symAttrQueryMap(a) = \lex$
		\ElsIf{$\symRMLGraph$ contains a triple $(u, \ttlSmaller{rr:constant}, o)$}
    		\State $\extExpr \gets o$
		\ElsIf{$\symRMLGraph$ contains a triple $(u, \ttlSmaller{rr:template}, o)$
			%such that
			s.t.\
		$o$ is a literal $\literalTuple$}
            \State  $\extExpr \gets $ \textsc{Template}($\lex, \symAttrQueryMap$)  \Comment Algorithm~\ref{alg:template_function}
		\EndIf
\medskip
        \If{$\symRMLGraph$ contains the triple $(u,\ttlSmaller{rr:termType}, \ttlSmaller{rr:BlankNode})$}
            \State \Return $\fctToBNode{\extExpr}$
        \ElsIf{$\symRMLGraph$ contains the triple $(u,\ttlSmaller{rr:termType}, \ttlSmaller{rr:IRI})$}
            \State \Return $\fctToIRI{\extExpr}$
        \ElsIf{$\symRMLGraph$ contains a triple $(u, \ttlSmaller{rr:datatype}, o)$ such that $o \in \symAllIRIs$}
            \State \Return $\fctToLiteral{\extExpr, o}$
        \Else 
            \State \Return $\fctToLiteral{\extExpr, \ttlSmaller{xsd:string}}$
		\EndIf
	\end{algorithmic}
\end{algorithm}

\end{algorithm}

The extend expression
	%that Algorithm~\ref{alg:rml_termmap_extend_expr} creates for a given term map depends on whether the
	%that Algorithm~\ref{alg:rml_termmap_extend_expr} creates depends on whether the given
	created by Algorithm~\ref{alg:rml_termmap_extend_expr} depends on whether the given
term map is constant-val\-ued
	%(lines~\ref{line:rml_termmap_extend_expr:from_constant_begin}--\ref{line:rml_termmap_extend_expr:from_constant_end}%
	(line~\ref{line:rml_termmap_extend_expr:from_constant_begin}%
%
	%\ of Algorithm~\ref{alg:rml_termmap_extend_expr}%
), reference-val\-ued (lines~\ref{line:rml_termmap_extend_expr:from_reference_begin}--\ref{line:rml_termmap_extend_expr:from_reference_end}), or template-val\-ued (lines~\ref{line:rml_termmap_extend_expr:from_template_begin}--\ref{line:rml_termmap_extend_expr:from_template_end}). In the latter case, the extend expression needs to
	%be created such that its evaluation instantiates
	be created such that it instantiates
	%instantiate
the corresponding template string. To this end, this string is first split into a sequence of substrings (line~\ref{line:rml_termmap_extend_expr:split_template}), defined as follows.
\begin{definition}
	\label{def:tempSubStrs}
	\normalfont
	For every string $\mathit{tmpl} \in \symAllStrings$,
	    %(assumed to be an RML \emph{string template}).
	we write $\tempSubStrs(\mathit{tmpl})$ to denote a sequence~$\substrSeq$ of strings that is constructed by the following
		two %four
	steps.
	\begin{enumerate}
		\item
		Partition $\mathit{tmpl}$ into a sequence~$\substrSeq$ of \emph{query substrings} and \emph{normal substrings} in the order in which these substrings appear in $\mathit{tmpl}$ from left to right,~where
		\begin{itemize}
			\item a \emph{query substring} is a substring
				starting with an opening curly brace (i.e., "\{") and ending with a closing curly brace (i.e., "\}"), and
				%starting and ending with curly braces (i.e., \{ and \}), and
				%enclosed in curly braces (i.e., \{ and \}), and
			\item a \emph{normal substring} is every substring between two query substrings, as well as
				%the substring before the first query substring and the substring
				the one before the first and the one
				%before the first and
			after the last query substring.
		\end{itemize}

		\item 
		If $\substrSeq$ is the empty sequence, insert the empty string as a new substring into~$\substrSeq$.

% 		\item
% 		For every two query substrings $q_1$ and $q_2$ that occur directly after one another in $\substrSeq$, insert the empty string as a new substring between $q_1$ and $q_2$ in $\substrSeq$.
% 
% 		\item
% 		If the first element in $\substrSeq$ is a query substring $q$, insert the empty string as a new substring before $q$ in $\substrSeq$. Similarly, if the last element in $\substrSeq$ is a query substring $q'\!$, insert the empty string as a new substring after $q'$ in $\substrSeq$.
	\end{enumerate}
\end{definition}

% 	%By the steps in Definition~\ref{def:tempSubStrs}, the
% 	By Definition~\ref{def:tempSubStrs}, the
% 	%The
% %
% 	%resulting sequence
% 	%sequence~$\substrSeq$ that Algorithm~\ref{alg:rml_termmap_extend_expr} uses
% 	sequence~$\substrSeq$ in Algorithm~\ref{alg:rml_termmap_extend_expr}
% %
% 	%is not empty, has an odd number of elements,
% 	has an odd number of elements
% and alternates between normal substrings and query substrings, beginning with a normal substring (which may be the empty string).
%
The extend expression is then created to concatenate the
	resulting
substrings~(lines~\ref{line:rml_termmap_extend_expr:first_substring}--\ref{line:rml_termmap_extend_expr:concatenation_end}), where every query substring "\{$q$\}" is replaced by the attribute~$\attr$ for which
	it holds that
$\symAttrQueryMap(\attr) = q$, because that attribute holds the values that the source operator obtains for the query
	expression~%
$q$.

\begin{example}
	\label{ex:rml_termmap_extend_expr}
	Assume
		%$\symRMLGraph_\mathsf{ex}$ is the RDF graph in Listing~\ref{lst:rml_example}, $b_\mathsf{sm}$ and $b_\mathsf{om}$ are the blank nodes that denote the subject map and the object map in this graph, respectively,
		$b_\mathsf{sm}$ and $b_\mathsf{om}$ are the blank nodes denoting the subject map and the object map in RDF graph~$\symRMLGraph_\mathsf{ex}$ in Listing~\ref{lst:rml_example}, respectively,
	and
		$\symAttrQueryMap = \{ \attr_1 \rightarrow \textsf{\small ID},$ $\attr_2 \rightarrow \textsf{\small Name} \}$
		%$\symAttrQueryMap$
	as in Example~\ref{ex:ExtractQueries}.
	Then, $\textsc{CreateExtExpr}(b_\mathsf{sm}, \symRMLGraph_\mathsf{ex}, \LTB, \symAttrQueryMap)$
		%returns the extend expression
		%is the extend expression
		returns
	$
	\toIRI\bigl(
		\concat\bigl(
				%(\ttl{"http://example.com/person\_"}, \ttl{xsd:string}),
				\ell,
			\attr_1
		\bigr)
	\bigr)
	$
		where $\ell$ is the literal $(\ttl{"http://example.com/person\_"},$ $\ttl{xsd:string})$%
	%
		%, and
		;
	$\textsc{CreateExtExpr}(b_\mathsf{om}, \symRMLGraph_\mathsf{ex}, \LTB, \symAttrQueryMap)$
		%returns
		is
	$\fctToLiteral{\attr_2, \ttl{xsd:string}}$.
\end{example}


\smallskip\noindent\textbf{\itshape Translation of the Object Map.}
For object (term) maps, the algorithm also has to consider the case of
	%so-called \emph{referencing object maps}%
	referencing object maps%
~(lines~\ref{line:rml2algebra:join_case_begin}--\ref{line:rml2algebra:join_case_end} in Algorithm~\ref{alg:rml2algebra}), which capture a join with
	RDF terms defined by the subject map
	%subjects
of another triples map. In this case, another source operator needs to be added
	(for this other triples map)
and
	the mapping relation created by that source operator needs
	%its mapping relation
be joined with the mapping relation produced before (line~\ref{line:rml2algebra:add_equijoin}). To this end, an equijoin operator is used, for which the join attributes are determined based on the join condition defined for the referencing object map~(lines~\ref{line:rml2algebra:init_join_pairs}--\ref{line:rml2algebra:populate_join_pairs_end}).
%
Ordinary object maps are handled
	%in the same way as
	like
subject
	%maps
and predicate maps before~(line~\ref{line:rml2algebra:extend_for_ordinary_object}).


\smallskip\noindent\textbf{\itshape Translation of the Graph Map.}
Finally, if the triples map has a graph map, the algorithm adds an extend operator, with attribute~$\graphAttr$, to create IRIs defined by this graph map (lines~\ref{line:rml2algebra:extend_for_graph_map_begin}--\ref{line:rml2algebra:extend_for_graph_map_end}). If the triples map
	%does not have a
	has no
graph map, an extend operator that assigns $\graphAttr$ to the constant IRI \ttl{rr:defaultGraph} is added~(line~\ref{line:rml2algebra:default_graph}).

